\chapter{非惯性系中的运动}

\section{旋转基矢的时间导数}
设移动坐标系的原点在惯性系中的位置为$\bm R(t)$，其角速度为$\bm\Omega(t)$。
对任意矢量$\bm A$，同一时刻在同一物理空间比较导数有
\begin{equation}
 \Bigp{\dv{\bm A}{t}}_\inertial
 = \Bigp{\dv{\bm A}{t}}_\rotating + \bm\Omega \times \bm A.
 \label{eq:cm-transport}
\end{equation}
第二项来自基矢的转动。将$\bm r_\inertial = \bm R + \bm r$连续求导两次得到
\begin{equation}
 \bm a_\inertial = \ddot{\bm R} + \bm a_\rotating
 + 2\bm\Omega \times \bm v_\rotating
 + \dotbm{\Omega} \times \bm r
 + \bm\Omega \times (\bm\Omega \times \bm r).
\end{equation}
因此旋转系中的运动方程是
\begin{align}
 m\bm a_\rotating = {}&\bm F - m\ddot{\bm R}
 - 2m\bm\Omega \times \bm v_\rotating
 - m\dotbm{\Omega} \times \bm r\notag\\
 &-m\bm\Omega \times (\bm\Omega \times \bm r).
 \label{eq:cm-rotating-newton}
\end{align}
右边依次为真实力、平移惯性力、Coriolis力、Euler力和离心力。
它们是变换方程产生的项，不对应一个另外施力的物体。
对恒定$\bm\Omega$，离心项向外，Coriolis项与旋转系速度垂直而不做功。

\section{旋转系中的能量}
若原点不加速、$\bm\Omega$恒定，且真实力为$\bm F = -\grad V$，其中$V$在旋转坐标中不显含时间，
则离心力可写成$-\grad V_{\mathrm{cf}}$，其中
\begin{equation}
 V_{\mathrm{cf}} = -\frac12m\abs{\bm\Omega \times \bm r}^2.
\end{equation}
此时
\begin{equation}
 E_\rotating = \frac12m v_\rotating^2 + V(\bm r)
 - \frac12m\abs{\bm\Omega \times \bm r}^2
\end{equation}
守恒。该量不等于惯性系机械能；旋转驱动装置可以向物体输入或取走能量。
若$V$在旋转系中仍显含时间，或$\bm\Omega$随时间变化，上述守恒式通常失效。

\begin{example}{Foucault摆的进动方向}{cm-foucault}
在纬度$\lambda$处，取$x$向东、$y$向北、$z$竖直向上。
对小振幅、长度远小于地球半径的摆，将离心修正吸收到当地有效$g$中，
保留地球自转角速度的一阶项。记$\omega_0^2 = g/l$，地球自转角速度的竖直分量为$\Omega_z = \Omega_\oplus\sin\lambda$。忽略高阶竖直位移后，水平运动满足
\begin{subequations}
 \begin{align}
  \ddot x + \omega_0^2x &= 2\Omega_z\dot y,\\
  \ddot y + \omega_0^2y &= -2\Omega_z\dot x.
 \end{align}
\end{subequations}
令$z_{\mathrm c} = x + \i y$，则$\ddot z_{\mathrm c} + 2\i\Omega_z\dot z_{\mathrm c} + \omega_0^2z_{\mathrm c} = 0$。
代入$z_{\mathrm c} = \e^{-\i\Omega_z t}w$可得$\ddot w + (\omega_0^2 + \Omega_z^2)w = 0$。
由于原方程只保留自转的一阶效应，不据此解释二阶频率修正；一阶结果是摆动平面以
$-\Omega_\oplus\sin\lambda$进动，即北半球从上方看为顺时针。
赤道处一阶进动为零，极点处每个恒星日相对地面转一周。
\end{example}

\paragraph{复习与阅读}
先写清原点运动与转轴方向，再使用式\eqref{eq:cm-transport}，可避免背诵惯性力时丢符号。
旋转系静止平衡要求$\bm a_\rotating = \bm v_\rotating = 0$，
并不意味着惯性系加速度为零。延伸阅读：Taylor或Morin的非惯性参考系章节。
